\section{Méthodologie}\label{sec:methodo}
\subsection{Application de référence}
L'application est un gestionnaire de tâches (to-do) exposé comme une API REST en JSON. Les données sont
conservées dans SQLite : une table \texttt{categories} (20 lignes) et une table \texttt{tasks} (identifiant, titre,
description, statut parmi \textit{todo}/\textit{doing}/\textit{done}, priorité de 1 à 3, catégorie, échéance), avec
des index sur le statut et sur la catégorie. Toutes les versions partent de la même base de 10\,000 tâches, générée
avec une graine fixe. Cinq fonctionnalités sont soumises à la charge :
\textbf{F1} créer une tâche (écriture) ;
\textbf{F2} lister les tâches avec pagination (lecture et comptage) ;
\textbf{F3} rechercher par mot-clé (\texttt{LIKE}) avec filtre de statut facultatif (parcours de la table) ;
\textbf{F4} lister les tâches d'une catégorie (jointure) ;
\textbf{F5} calculer des statistiques par statut et par catégorie (agrégation).
Trois routes complètent le CRUD (lire, modifier le statut, supprimer) ; elles sont vérifiées par le test
d'équivalence mais ne sont pas chargées. Le cahier des charges complet (routes, validations, codes d'erreur) est joint
au rendu.

\subsection{Technologies comparées}
Nous comparons cinq technologies (tableau~\ref{tab:techs}), qui couvrent plusieurs modèles d'exécution : code natif
compilé (Rust, Go), JavaScript compilé à la volée sur une boucle d'événements (Node.js) et langages interprétés
(PHP, Python). PHP et Python sont aussi des langages très répandus pour le développement web, ce qui rend
la comparaison utile aux étudiants. Chaque version utilise le serveur HTTP le plus simple disponible dans sa
technologie, sans paquet externe (voir la section~\ref{sec:dev}).

\begin{table}[t]
\caption{Technologies comparées}\label{tab:techs}
\centering\small
\begin{tabular}{@{}llll@{}}
\toprule
\textbf{Technologie} & \textbf{Langage} & \textbf{Version} & \textbf{SQLite}\\
\midrule
Flask 3.1.3 & Python & 3.13.18 & 3.53.1\\
module \texttt{http} & Node.js & 22.23.4 & 3.53.4\\
\texttt{php -S} & PHP & 8.3.6 & 3.45.1\\
\texttt{net/http} & Go & 1.24.7 & 3.45.1\\
bibl. standard & Rust & 1.97.0 & 3.45.1\\
\bottomrule
\end{tabular}
\end{table}

\subsection{Métriques}
Une seule grandeur ne suffit pas pour dire qu'une version est « plus performante » : un serveur peut répondre vite
en consommant beaucoup de processeur, ou tenir la charge en répondant lentement à certains utilisateurs. Nous
retenons six familles de mesures, relevées pour chaque essai.
\begin{itemize}[leftmargin=*,itemsep=0pt,topsep=2pt]
\item \textbf{Débit} (requêtes/s) : nombre de requêtes traitées divisé par la durée de l'essai ; mesure la capacité
à absorber du trafic. Outil : générateur de charge.
\item \textbf{Latence moyenne, médiane (p50), p95 et p99} (ms) : durée entre l'envoi de la requête et la fin de la
réponse, vue par le client. La médiane décrit le cas courant, p95 et p99 l'expérience des utilisateurs les moins
bien servis, que la moyenne masque \cite{dean2013}.
\item \textbf{Taux d'erreur} : requêtes échouées (réseau ou code HTTP hors 2xx) ; un débit n'a de sens que sans erreur.
\item \textbf{Temps CPU par requête} (ms) : temps processeur (utilisateur + système) du serveur pendant l'essai,
divisé par le nombre de requêtes ; mesure le coût de chaque requête, donc l'efficacité énergétique et financière.
Outil : \texttt{psutil}, d'après \texttt{/proc}.
\item \textbf{Mémoire} (Mo) : \textit{proportional set size} (PSS) de tous les processus du serveur en fin d'essai
(plus fidèle que la RSS lorsque plusieurs processus partagent des pages, comme PHP).
\item \textbf{Temps de démarrage} (ms) : du lancement du processus à la première réponse correcte ; utile pour
le déploiement et la mise à l'échelle.
\end{itemize}

\subsection{Outils et environnement de test}
Le générateur de charge est un programme Go écrit pour l'occasion (aucun outil comme k6 ou wrk n'était disponible
sur la machine, sans accès réseau). Chaque « utilisateur simultané » est une goroutine qui garde sa connexion HTTP/1.1
ouverte et envoie la requête suivante dès la réponse reçue (boucle fermée, sans temps de réflexion \cite{jain1991}). Les tests
s'exécutent dans un conteneur Linux (Ubuntu 24.04, noyau 6.18) à 2 cœurs virtuels (Xeon à 2,1\,GHz) et 8\,Go de
mémoire. Le serveur est épinglé sur le cœur~0 et le générateur sur le cœur~1 (\texttt{taskset}), si bien que
\emph{chaque version dispose d'un seul cœur} : les résultats mesurent le rendement par cœur, non l'avantage du
parallélisme.

\subsection{Déroulement des tests}
Pour chaque fonctionnalité (F1 à F5), chaque niveau de charge (10, 50 et 100 utilisateurs simultanés) et chaque
essai numéroté de 1 à 30, nous exécutons, pour les cinq technologies dans un ordre tiré au sort (graine fixe) :
(1)~copie d'une base neuve identique et démarrage du serveur ; (2)~100 requêtes de chauffe, non comptées ;
(3)~essai mesuré de 1\,000 requêtes (500 pour F3 et F5, plus longues) réparties entre les utilisateurs ;
(4)~relevé du CPU et de la mémoire, puis arrêt du serveur. L'ordre aléatoire répartit sur toutes les technologies
les dérives de la machine (cache, voisins du conteneur) ; le redémarrage à chaque essai rend les essais
indépendants, y compris pour F1 qui modifie la base. Les 30 essais par cas permettent de calculer des moyennes
avec intervalles de confiance à 95\,\% et des tests de comparaison \cite{georges2007}.
Les scripts, le générateur et les données brutes sont fournis pour répéter l'expérience.
